\documentclass{ximera}




% MATH -----------------------------------------------------------
\newcommand{\nats}{\mathbb N}
\newcommand{\ints}{\mathbb Z}
\newcommand{\rats}{\mathbb Q}
\newcommand{\reals}{\mathbb R}
\newcommand{\complex}{\mathbb C}
\newcommand{\powerset}{\mathscr P}

%-------------------------------------------------------------

\pgfplotsset{compat=1.18}
\begin{document}
\begin{abstract}
 \end{abstract}

    \title{Class Discussion Activity 22}
    
    \maketitle

\section*{Exercises}

\begin{enumerate}[label=(\arabic*)]


\item Suppose $f(t)$ is a continuous function on $[0,\infty)$ and $\mathcal{L}(f(t))=\frac{5}{s^2+1}$.  Which is the Laplace transform of the solution to $y\,^{\prime}-2y=f(t)$, $y(0)=0$?

\begin{enumerate}[label=(\alph*)]
\item $\dfrac{5s-10}{s^2+1}$
\item $\dfrac{5}{s^3-s^2-2}$
\item $\dfrac{5}{s^3-2s^2+s-2}$  % ***
\item $\dfrac{5}{s^3-s^2+2}$
\item $\dfrac{5}{s^3-2s^2-s+2}$ 
\item $\dfrac{5}{s^3-s^2-2s-2}$ 
\end{enumerate}

% (c) with options through (f)

\item Which is the Laplace transform of the solution to the IVP \\$y^{\,\prime\prime}-y^{\,\prime}-2y=27te^{-t}$, $y(0)=0$, $y^{\,\prime}(0)=0$?

\begin{enumerate}[label=(\alph*)]
\item $\dfrac{27}{(s+1)^2}$
\item $\dfrac{27}{(s+1)^2(s-2)}$
\item $\dfrac{27s}{(s+1)^2(s-2)}$
\item $\dfrac{27}{(s+1)^3(s-2)}$  % ***
\item $\dfrac{27s}{(s+1)^3(s-2)}$ 
\item $\dfrac{27s^2}{(s+1)^3(s-2)}$
\end{enumerate}

% (d) with options through (f)


\newpage

\item Determine $y(1)$ given that $y(t)$ satisfies the IVP in the previous exercise.  Round to the nearest hundredth.

% 4.26


\item Suppose that $f(t)$ is a differentiable function on $[0,\infty)$ and that $f^{\,\prime}(t)$ is exponential order $s_0$ for some real number $s_0$.  Determine $$\displaystyle \lim_{s\rightarrow\infty} s\mathcal{L}(f(t))-f(0).$$

Hints: $\displaystyle s\mathcal{L}(f(t))-f(0)=\mathcal{L}(f\,^{\prime}(t))=\int_0^{\infty} f^{\,\prime}(t)e^{-st}\,\mathrm{d}t$, and since $f^{\,\prime}(t)$ is of exponential order $s_0$, the limit as $s\rightarrow \infty$ of the improper integral is equal to the improper integral of the limit of the integrand.

% 0



\item Let $a,b,c$ be real numbers and assume $a\neq 0$. Let $y(t)=\mathcal{L}^{-1}\left(\dfrac{as+b}{as^2+bs+c}\right)$.  Determine $y(0)$.  

Hint: use the result of the previous exercise.


\begin{enumerate}[label=(\alph*)]
\item $0$
\item $1$ % ***
\item $2$
\item $a$
\item $1/a$
\item $b$
\item $b/a$
\item None of these
\end{enumerate}

% (b) with options through (h)

\end{enumerate}

\end{document}